\documentclass[handout,12pt,aspectratio=169]{beamer}
\usepackage[T1]{fontenc}
\usepackage{lmodern}
\usepackage{amsmath,amssymb,array,booktabs,tabularx,tikz}
% Add language support appropriate to the source, if needed.
% Keep this preamble inline in every delivered deck.
\definecolor{GEInk}{HTML}{132A3E}
\definecolor{GETeal}{HTML}{17675E}
\definecolor{GEFooter}{HTML}{EAF2F1}
\setbeamercolor{background canvas}{bg=white}
\setbeamercolor{normal text}{fg=GEInk,bg=white}
\setbeamercolor{structure}{fg=GETeal}
\setbeamercolor{frametitle}{fg=GEInk}
\setbeamercolor{GE footer}{fg=GETeal,bg=GEFooter}
\setbeamerfont{normal text}{family=\sffamily,size=\normalsize}
\setbeamerfont{title}{size=\huge,series=\bfseries}
\setbeamerfont{subtitle}{size=\normalsize,series=\bfseries}
\setbeamerfont{frametitle}{size=\Large,series=\bfseries}
\setbeamerfont{framesubtitle}{size=\normalsize,series=\mdseries}
\setbeamerfont{GE footer}{size=\footnotesize}
\setbeamersize{text margin left=9mm,text margin right=9mm}
\setbeamertemplate{navigation symbols}{}
\setbeamertemplate{headline}{}
\setbeamercovered{invisible}
\setbeamertemplate{frametitle}{%
  \vskip.5em
  {\usebeamerfont{frametitle}\strut\insertframetitle\strut\par}%
  \ifx\insertframesubtitle\empty\else
    {\usebeamerfont{framesubtitle}\strut\insertframesubtitle\strut\par}%
  \fi
  \vskip.25em
}
\setbeamertemplate{footline}{%
  \begin{beamercolorbox}[wd=\paperwidth,ht=3ex,dp=1.3ex,
    leftskip=9mm,rightskip=9mm]{GE footer}%
    \usebeamerfont{GE footer}\insertsectionhead\hfill
    \insertframenumber/\inserttotalframenumber
  \end{beamercolorbox}%
}
\setlength{\parindent}{0pt}
\setlength{\parskip}{.35em}
\setlength{\tabcolsep}{4pt}
\renewcommand{\arraystretch}{1.22}
\AtBeginEnvironment{frame}{\raggedright}
\addtobeamertemplate{frametitle}{}{%
  \setlength{\abovedisplayskip}{6pt plus 2pt minus 2pt}%
  \setlength{\belowdisplayskip}{6pt plus 2pt minus 2pt}%
  \setlength{\abovedisplayshortskip}{6pt plus 2pt minus 2pt}%
  \setlength{\belowdisplayshortskip}{6pt plus 2pt minus 2pt}%
}


% SOURCE MAP / user-directed sequence:
% Legend: practice.tex, document-local practicehighlightlegend@G11-T2-C2 override.
% Example: classroom_activity.tex, problem 1, narrative + data + Tasks 1--3.
% No guided/independent exercises. Legacy simulation labels normalized to
% mu_known: the same fixed population mean, later established as 41 minutes.
\definecolor{PopulationMean}{HTML}{FF765C}
\definecolor{SampleMean}{HTML}{9BDD36}
\definecolor{PopulationSD}{HTML}{E42134}
\definecolor{SampleSize}{HTML}{173B17}
\definecolor{Confidence}{HTML}{55C8F0}
\definecolor{Critical}{HTML}{2648BB}
\definecolor{StandardError}{HTML}{F4D94B}
\definecolor{MarginError}{HTML}{F28A18}
\colorlet{Endpoint}{MarginError!45!white}
\colorlet{Width}{MarginError!45!black}
\definecolor{Coverage}{HTML}{BC9DE0}

\definecolor{Error}{HTML}{28AFA2}
\newif\ifGTenHighlights
\GTenHighlightstrue
% \GTenHighlightsfalse % Uncomment to disable instructional highlights.
% Transparent markers retain identical dimensions before activation and with
% highlights disabled. Visibility and instructional activation are independent.
\tikzset{ge marker/.style={rectangle,rounded corners=1mm,draw=none,
 fill=#1,fill opacity=\GEOpacity,text opacity=1,text=GEInk,
 inner xsep=.6mm,inner ysep=.3mm,outer sep=0pt}}
\makeatletter
\newcommand{\GEHighlightMath}[2]{\tikz[baseline=(gehl.base)]
 \node[ge marker=\GECurrentRole](gehl){$\m@th#1#2$};}
\newcommand{\GEHighlight}[2]{\begingroup
 \def\GEOpacity{0}\ifGTenHighlights\def\GEOpacity{.32}\fi
 \def\GECurrentRole{#1}\ifmmode\mathpalette\GEHighlightMath{#2}
 \else\tikz[baseline=(gehl.base)]\node[ge marker=#1](gehl){#2};\fi\endgroup}
\makeatother
\newcommand{\GEPlain}[1]{{\GTenHighlightsfalse\GEHighlight{white}{#1}}}
\newcommand{\At}[3]{\ifnum#1=0\relax\GEPlain{#3}\else\alt<#1->{\GEHighlight{#2}{#3}}{\GEPlain{#3}}\fi}
\newcommand{\R}[2]{\leavevmode\visible<#1->{#2}}
\newcommand{\CtwoMean}[1]{\GEHighlight{PopulationMean}{#1}}
\newcommand{\CtwoPoint}[1]{\GEHighlight{SampleMean}{#1}}
\newcommand{\CtwoSD}[1]{\GEHighlight{PopulationSD}{#1}}
\newcommand{\CtwoSize}[1]{\GEHighlight{SampleSize}{#1}}
\newcommand{\CtwoConfidence}[1]{\GEHighlight{Confidence}{#1}}
\newcommand{\CtwoSignificance}[1]{\GEHighlight{Confidence}{#1}}
\newcommand{\CtwoCritical}[1]{\GEHighlight{Critical}{#1}}
\newcommand{\CtwoSE}[1]{\GEHighlight{StandardError}{#1}}
\newcommand{\CtwoMoE}[1]{\GEHighlight{MarginError}{#1}}
\newcommand{\CtwoEndpoint}[1]{\GEHighlight{Endpoint}{#1}}
\newcommand{\CtwoCI}[1]{\GEHighlight{Endpoint}{#1}}
\newcommand{\CtwoWidth}[1]{\GEHighlight{Width}{#1}}
\newcommand{\CtwoCoverage}[1]{\GEHighlight{Coverage}{#1}}
\newcommand{\CtwoError}[1]{\GEHighlight{Error}{#1}}
\newcommand{\CtwoCoverCount}{\CtwoCoverage{N_{\rm cover}}}
\newcommand{\CtwoTotalCount}{\CtwoCoverage{N}}
\newcommand{\LegendStyle}{\fontsize{9.5}{10}\selectfont\setlength{\parskip}{0pt}\setlength{\tabcolsep}{2pt}\renewcommand{\arraystretch}{1.0}\setlength{\aboverulesep}{1pt}\setlength{\belowrulesep}{1pt}\tikzset{ge marker/.append style={inner xsep=.4mm,inner ysep=.1mm}}}
\newcommand{\LeftLegend}{\begingroup\LegendStyle\begin{tabularx}{\linewidth}{@{}>{\raggedright\arraybackslash}p{.28\linewidth}>{\raggedright\arraybackslash}p{.30\linewidth}>{\raggedright\arraybackslash}X@{}}
\toprule
\textbf{Symbol} & \textbf{Concept} & \textbf{Formula} \\
\midrule
$\CtwoConfidence{C}$ & \CtwoConfidence{Confidence} \CtwoConfidence{level} &  \\[1pt]
$\CtwoCritical{z^*}$ & \CtwoCritical{Critical} \CtwoCritical{value} &  \\[1pt]
$\CtwoSE{\mathrm{SE}}$ & \CtwoSE{Standard} \CtwoSE{error} & $\CtwoSE{\mathrm{SE}}=\dfrac{\CtwoSD{\sigma}}{\sqrt{\CtwoSize{n}}}$ \\[1pt]
$\CtwoMoE{\mathrm{MoE}}$ & \CtwoMoE{Margin} \CtwoMoE{of} \CtwoMoE{error} & $\CtwoMoE{\mathrm{MoE}}=\CtwoCritical{z^*}\times\CtwoSE{\mathrm{SE}}$ \\[1pt]
$\CtwoEndpoint{L}$ & \CtwoEndpoint{Lower} \CtwoEndpoint{endpoint} & $\CtwoEndpoint{L}=\CtwoPoint{\bar{x}}-\CtwoMoE{\mathrm{MoE}}$ \\[1pt]
$\CtwoEndpoint{U}$ & \CtwoEndpoint{Upper} \CtwoEndpoint{endpoint} & $\CtwoEndpoint{U}=\CtwoPoint{\bar{x}}+\CtwoMoE{\mathrm{MoE}}$ \\[1pt]
\bottomrule
\end{tabularx}\par\endgroup}
\newcommand{\RightLegend}{\begingroup\LegendStyle\begin{tabularx}{\linewidth}{@{}>{\raggedright\arraybackslash}p{.28\linewidth}>{\raggedright\arraybackslash}p{.30\linewidth}>{\raggedright\arraybackslash}X@{}}
\toprule
\textbf{Symbol} & \textbf{Concept} & \textbf{Formula} \\
\midrule
$\CtwoCI{\mathrm{CI}}$ & \CtwoCI{Confidence} \CtwoCI{interval} & $\CtwoCI{\mathrm{CI}}=[\CtwoEndpoint{L},\CtwoEndpoint{U}]$ \\[1pt]
$\CtwoWidth{W}$ & \CtwoWidth{Full} \CtwoWidth{width} & $\CtwoWidth{W}=\CtwoEndpoint{U}-\CtwoEndpoint{L}$\par$=2\CtwoMoE{\mathrm{MoE}}$ \\[1pt]
$\CtwoSignificance{\alpha}$ & \CtwoSignificance{Significance} \CtwoSignificance{level} & $\CtwoSignificance{\alpha}=1-\CtwoConfidence{C}$ \\[1pt]
$\CtwoSignificance{\alpha}/2$ & \CtwoSignificance{Each} \CtwoSignificance{tail}'s area & $\CtwoSignificance{\alpha}/2=\dfrac{1-\CtwoConfidence{C}}{2}$ \\[1pt]
$\CtwoCoverCount(\CtwoConfidence{C})$ & \CtwoCoverage{Covering} \CtwoCoverage{count} &  \\[1pt]
$\CtwoTotalCount(\CtwoConfidence{C})$ & \CtwoCoverage{Total} \CtwoCoverage{interval} \CtwoCoverage{count} &  \\[1pt]
$\CtwoCoverage{\mathrm{Cov}}(\CtwoConfidence{C})$ & \CtwoCoverage{Observed} \CtwoCoverage{coverage} & $\dfrac{\CtwoCoverCount(\CtwoConfidence{C})}{\CtwoTotalCount(\CtwoConfidence{C})}\times100\%$ \\[1pt]
\bottomrule
\end{tabularx}\par\endgroup}
\newcommand{\BothLegends}{\begin{columns}[T,onlytextwidth]
\begin{column}{.49\textwidth}\LeftLegend\end{column}
\begin{column}{.49\textwidth}\RightLegend\end{column}\end{columns}}

\newif\ifReference
\Referencefalse
% In reference mode all entries are visible; each frame supplies input activation.
\newcommand{\IR}[6]{\ifReference
 \At{\csname Input#1\endcsname}{#2}{$#3$}&\Words{\csname Input#1\endcsname}{#2}{#4}&\Words{\csname Input#1\endcsname}{#2}{#5}\\[1pt]
 \else
 \At{#1}{#2}{$#3$}&\Words{#1}{#2}{#4}&\R{#1}{\Words{#1}{#2}{#5}}\\[1pt]
 \fi}
\newcommand{\InfoTables}{\begingroup
\fontsize{9.5}{10}\selectfont\setlength{\parskip}{0pt}\setlength{\tabcolsep}{2pt}\renewcommand{\arraystretch}{1.0}\setlength{\aboverulesep}{.7pt}\setlength{\belowrulesep}{.7pt}\tikzset{ge marker/.append style={inner xsep=.25mm,inner ysep=.1mm}}
\begin{columns}[T,onlytextwidth]
\begin{column}{.63\textwidth}\begin{tabularx}{\linewidth}{@{}>{\raggedright\arraybackslash}p{.23\linewidth}>{\raggedright\arraybackslash}p{.38\linewidth}>{\raggedright\arraybackslash}X@{}}\toprule
\textbf{Symbol}&\textbf{Concept}&\textbf{In context...}\\\midrule
\multicolumn{3}{@{}l}{\textbf{Population parameters}}\\\midrule
\IR{2}{PopulationMean}{\mu}{Population mean}{Month's mean trip duration (minutes)}{}
\IR{3}{PopulationMean}{\mu_{\mathrm{known}}}{Later established population mean}{41 minutes}{}
\IR{4}{PopulationSD}{\sigma}{Population standard deviation}{14 minutes}{}
\multicolumn{3}{@{}l}{\textbf{Sample statistics}}\\\midrule
\IR{5}{SampleMean}{\bar{x}}{Sample mean}{40 minutes}{}
\IR{6}{SampleSize}{n}{Sample size}{49 trips}{}
\IR{7}{StandardError}{\mathrm{SE}}{Standard error}{2 minutes}{}
\end{tabularx}\end{column}
\begin{column}{.35\textwidth}\begin{tabularx}{\linewidth}{@{}>{\raggedright\arraybackslash}p{.23\linewidth}>{\raggedright\arraybackslash}p{.38\linewidth}>{\raggedright\arraybackslash}X@{}}\toprule
\textbf{Symbol}&\textbf{Concept}&\textbf{In context...}\\\midrule
\multicolumn{3}{@{}l}{\textbf{Procedure Settings}}\\\midrule
\IR{9}{Confidence}{C}{Confidence level}{95\% = 0.95}{}
\IR{11}{Confidence}{\alpha}{Significance level}{0.05}{}
\IR{13}{Confidence}{\alpha/2}{Tail area}{0.025}{}
\IR{15}{Critical}{z^*}{Critical value}{1.960}{}
\end{tabularx}\par\smallskip\ifReference\else \begin{overlayarea}{\linewidth}{1.15cm}\fontsize{9.5}{10}\selectfont
\only<7>{$\CtwoSE{\mathrm{SE}}=\dfrac{\CtwoSD{14}}{\sqrt{\CtwoSize{49}}}=\CtwoSE{2}$ minutes.}
\only<9>{$\CtwoConfidence{C}=\CtwoConfidence{95\%}=\CtwoConfidence{0.95}$.}
\only<11>{$\CtwoSignificance{\alpha}=1-\CtwoConfidence{0.95}=\CtwoSignificance{0.05}$.}
\only<13>{$\CtwoSignificance{\alpha}/2=\CtwoSignificance{0.05}/2=\CtwoSignificance{0.025}$.}
\only<15>{$\CtwoCritical{z^*}=\CtwoCritical{1.960}$: standard normal critical value for two-sided \CtwoConfidence{95\%} confidence.}
\end{overlayarea}\fi\end{column}
\end{columns}\endgroup}
\ExplSyntaxOn
\NewDocumentCommand{\Words}{mmm}{\seq_set_split:Nnn \l_tmpa_seq {~}{#3}\seq_map_inline:Nn \l_tmpa_seq {\At{#1}{#2}{##1}\space}}
\ExplSyntaxOff
\expandafter\def\csname Input2\endcsname{0}
\expandafter\def\csname Input3\endcsname{0}
\expandafter\def\csname Input4\endcsname{0}
\expandafter\def\csname Input5\endcsname{0}
\expandafter\def\csname Input6\endcsname{0}
\expandafter\def\csname Input7\endcsname{0}
\expandafter\def\csname Input9\endcsname{0}
\expandafter\def\csname Input11\endcsname{0}
\expandafter\def\csname Input13\endcsname{0}
\expandafter\def\csname Input15\endcsname{0}
\title{Point estimates and confidence levels}
\subtitle{Grade 11 · Term 2 · Lesson 1\\Global Economics · 2026--2027}
\author{}\date{}
\newcommand{\GEOfficialCriterion}{\textbf{C2} Compares point and interval estimation and interprets confidence levels through their long-run coverage meaning, distinguishing this from the probability that a particular calculated interval contains the parameter.}
\setbeamertemplate{title page}{\raggedright\vfill
{\color{GETeal}\usebeamerfont{subtitle}\insertsubtitle\par}\medskip
{\color{GEInk}\usebeamerfont{title}\inserttitle\par}\bigskip
{\normalsize\GEOfficialCriterion\par}\vfill}
\begin{document}
\section{Highlight guide}\begin{frame}\titlepage\end{frame}
\begin{frame}[t]{Key terms and highlight guide -- left table}
\begin{minipage}{.76\linewidth}\LeftLegend\end{minipage}
\end{frame}
\begin{frame}[t]{Key terms and highlight guide -- right table}
\begin{minipage}{.76\linewidth}\RightLegend\end{minipage}
\end{frame}
\begin{frame}[t]{Key terms and highlight guide}
\BothLegends
\end{frame}
\section{Cycling study}
\begin{frame}[t]{Investigate the cycling interval}
\small An initial study estimated the then unknown average duration of all cycling trips in a specified month. A random sample of 49 trips averaged 40 minutes; the known population standard deviation was 14 minutes. The report used 95\% confidence.\par\medskip
\begin{columns}[T,onlytextwidth]\begin{column}{.56\textwidth}Five subsequent independent studies sampled the same month's population, each using 49 trips, the same standard deviation and 95\% procedure. A later comprehensive population investigation established an average of 41 minutes, enabling retrospective interval assessment.\end{column}\begin{column}{.42\textwidth}\begin{tabular}{@{}rrrr@{}}\toprule Study & $\bar{x}$ (min) & $L$ (min) & $U$ (min)\\\midrule
1&40&36.08&43.92\\
2&42&38.08&45.92\\
3&45&41.08&48.92\\
4&38&34.08&41.92\\
5&41&37.08&44.92\\
\bottomrule\end{tabular}\end{column}\end{columns}
\end{frame}
\begin{frame}[t]{Task 1 -- Link the narrative to the information}
\small
\fontsize{10}{11}\selectfont\setlength{\parskip}{0pt}\begin{minipage}[t][1.2cm][t]{\linewidth}
An initial study estimated the then unknown \At{2}{PopulationMean}{average duration} of \At{2}{PopulationMean}{all cycling trips} in a \At{2}{PopulationMean}{specified month}.
A random sample of \At{6}{SampleSize}{49 trips} averaged \At{5}{SampleMean}{40 minutes}; the known \At{4}{PopulationSD}{population standard deviation} was \At{4}{PopulationSD}{14 minutes}. The report used \At{9}{Confidence}{95\% confidence}.
\end{minipage}\par
\begin{minipage}[t][.45cm][t]{\linewidth}
\R{3}{A later comprehensive population investigation established an average of \At{3}{PopulationMean}{41 minutes}.}
\end{minipage}
\InfoTables

\onslide<16>

\end{frame}
\section{Interval solution}
\begin{frame}[t]{Task 2 -- Point estimate}
\Referencetrue
\expandafter\def\csname Input2\endcsname{0}
\expandafter\def\csname Input3\endcsname{0}
\expandafter\def\csname Input4\endcsname{0}
\expandafter\def\csname Input5\endcsname{2}
\expandafter\def\csname Input6\endcsname{0}
\expandafter\def\csname Input7\endcsname{4}
\expandafter\def\csname Input9\endcsname{0}
\expandafter\def\csname Input11\endcsname{0}
\expandafter\def\csname Input13\endcsname{0}
\expandafter\def\csname Input15\endcsname{4}
\InfoTables\smallskip\fontsize{10.5}{11.5}\selectfont\setlength{\parskip}{0pt}\setlength{\aboverulesep}{1pt}\setlength{\belowrulesep}{1pt}
\begin{tabularx}{\linewidth}{@{}p{.10\linewidth}p{.18\linewidth}X@{}}\toprule\textbf{Symbol}&\textbf{Concept}&\textbf{Calculation, Result and meaning in context}\\\midrule\At{2}{SampleMean}{$\bar{x}$}&\Words{2}{SampleMean}{Point estimate}&\R{2}{$\CtwoPoint{\bar{x}}=\CtwoPoint{40}$ minutes.}\par\R{3}{One estimate of the mean cycling-trip duration.}\\\bottomrule\end{tabularx}\end{frame}
\begin{frame}[t]{Task 2 -- Margin of error}
\Referencetrue
\expandafter\def\csname Input2\endcsname{0}
\expandafter\def\csname Input3\endcsname{0}
\expandafter\def\csname Input4\endcsname{0}
\expandafter\def\csname Input5\endcsname{1}
\expandafter\def\csname Input6\endcsname{0}
\expandafter\def\csname Input7\endcsname{4}
\expandafter\def\csname Input9\endcsname{0}
\expandafter\def\csname Input11\endcsname{0}
\expandafter\def\csname Input13\endcsname{0}
\expandafter\def\csname Input15\endcsname{4}
\InfoTables\smallskip\fontsize{10.5}{11.5}\selectfont\setlength{\parskip}{0pt}\setlength{\aboverulesep}{1pt}\setlength{\belowrulesep}{1pt}
\begin{tabularx}{\linewidth}{@{}p{.10\linewidth}p{.18\linewidth}X@{}}\toprule\textbf{Symbol}&\textbf{Concept}&\textbf{Calculation, Result and meaning in context}\\\midrule
\At{4}{MarginError}{$\mathrm{MoE}$}&\Words{4}{MarginError}{Margin of error}&\R{4}{$\CtwoMoE{\mathrm{MoE}}=\CtwoCritical{z^*}\times\CtwoSE{\mathrm{SE}}$}\par\R{5}{$=\CtwoCritical{1.960}\times\CtwoSE{2}$}\R{6}{$=\CtwoMoE{3.92}$ minutes.}\par\R{7}{Uncertainty in the mean duration: centre-to-endpoint distance.}\\\bottomrule\end{tabularx}\end{frame}
\begin{frame}[t]{Task 2 -- Build the confidence interval}
\Referencetrue
\expandafter\def\csname Input2\endcsname{0}
\expandafter\def\csname Input3\endcsname{0}
\expandafter\def\csname Input4\endcsname{0}
\expandafter\def\csname Input5\endcsname{1}
\expandafter\def\csname Input6\endcsname{0}
\expandafter\def\csname Input7\endcsname{1}
\expandafter\def\csname Input9\endcsname{0}
\expandafter\def\csname Input11\endcsname{0}
\expandafter\def\csname Input13\endcsname{0}
\expandafter\def\csname Input15\endcsname{1}
\InfoTables\smallskip\fontsize{10.5}{11.5}\selectfont\setlength{\parskip}{0pt}\setlength{\aboverulesep}{1pt}\setlength{\belowrulesep}{1pt}
\begin{tabularx}{\linewidth}{@{}p{.10\linewidth}p{.18\linewidth}X@{}}\toprule\textbf{Symbol}&\textbf{Concept}&\textbf{Calculation, Result and meaning in context}\\\midrule\At{2}{Endpoint}{$\mathrm{CI}$}&\Words{2}{Endpoint}{Confidence interval}&
\R{2}{$\CtwoEndpoint{L}=\CtwoPoint{\bar{x}}-\CtwoMoE{\mathrm{MoE}}$}\R{3}{$=\CtwoPoint{40}-\CtwoMoE{3.92}$}\R{4}{$=\CtwoEndpoint{36.08}$ min.}\par
\R{5}{$\CtwoEndpoint{U}=\CtwoPoint{\bar{x}}+\CtwoMoE{\mathrm{MoE}}$}\R{6}{$=\CtwoPoint{40}+\CtwoMoE{3.92}$}\R{7}{$=\CtwoEndpoint{43.92}$ min.}\par
\\\bottomrule\end{tabularx}\end{frame}
\begin{frame}[t]{Task 2 -- Assemble the confidence interval}
\Referencetrue
\expandafter\def\csname Input2\endcsname{0}
\expandafter\def\csname Input3\endcsname{0}
\expandafter\def\csname Input4\endcsname{0}
\expandafter\def\csname Input5\endcsname{1}
\expandafter\def\csname Input6\endcsname{0}
\expandafter\def\csname Input7\endcsname{1}
\expandafter\def\csname Input9\endcsname{0}
\expandafter\def\csname Input11\endcsname{0}
\expandafter\def\csname Input13\endcsname{0}
\expandafter\def\csname Input15\endcsname{1}
\InfoTables\smallskip\fontsize{10.5}{11.5}\selectfont\setlength{\parskip}{0pt}\setlength{\aboverulesep}{1pt}\setlength{\belowrulesep}{1pt}
\begin{tabularx}{\linewidth}{@{}p{.10\linewidth}p{.18\linewidth}X@{}}\toprule\textbf{Symbol}&\textbf{Concept}&\textbf{Calculation, Result and meaning in context}\\\midrule\At{2}{Endpoint}{$\mathrm{CI}$}&\Words{2}{Endpoint}{Confidence interval}&
\R{2}{$\CtwoCI{\mathrm{CI}}=[\CtwoEndpoint{L},\CtwoEndpoint{U}]$}\R{3}{$=[\CtwoEndpoint{36.08},\CtwoEndpoint{43.92}]$ minutes.}\par
\R{4}{Plausible values for the fixed population mean under the model.}\\\bottomrule\end{tabularx}\end{frame}
\begin{frame}[t]{Task 2 -- Full width}
\Referencetrue
\expandafter\def\csname Input2\endcsname{0}
\expandafter\def\csname Input3\endcsname{0}
\expandafter\def\csname Input4\endcsname{0}
\expandafter\def\csname Input5\endcsname{1}
\expandafter\def\csname Input6\endcsname{0}
\expandafter\def\csname Input7\endcsname{1}
\expandafter\def\csname Input9\endcsname{6}
\expandafter\def\csname Input11\endcsname{0}
\expandafter\def\csname Input13\endcsname{0}
\expandafter\def\csname Input15\endcsname{1}
\InfoTables\smallskip\fontsize{10.5}{11.5}\selectfont\setlength{\parskip}{0pt}\setlength{\aboverulesep}{1pt}\setlength{\belowrulesep}{1pt}
\begin{tabularx}{\linewidth}{@{}p{.10\linewidth}p{.18\linewidth}X@{}}\toprule\textbf{Symbol}&\textbf{Concept}&\textbf{Calculation, Result and meaning in context}\\\midrule\At{2}{Width}{$W$}&\Words{2}{Width}{Full width}&\R{2}{$\CtwoWidth{W}=\CtwoEndpoint{U}-\CtwoEndpoint{L}=2\CtwoMoE{\mathrm{MoE}}$}\par\R{3}{$=\CtwoEndpoint{43.92}-\CtwoEndpoint{36.08}=2(\CtwoMoE{3.92})$}\R{4}{$=\CtwoWidth{7.84}$ minutes.}\par\R{5}{The span of plausible mean durations describes the estimate's precision.}\\\bottomrule\end{tabularx}\par\smallskip\end{frame}
\begin{frame}[t]{Task 2 -- Interpret the interval}
\Referencetrue
\expandafter\def\csname Input2\endcsname{0}
\expandafter\def\csname Input3\endcsname{0}
\expandafter\def\csname Input4\endcsname{0}
\expandafter\def\csname Input5\endcsname{1}
\expandafter\def\csname Input6\endcsname{0}
\expandafter\def\csname Input7\endcsname{1}
\expandafter\def\csname Input9\endcsname{6}
\expandafter\def\csname Input11\endcsname{0}
\expandafter\def\csname Input13\endcsname{0}
\expandafter\def\csname Input15\endcsname{1}
\InfoTables\smallskip\fontsize{10.5}{11.5}\selectfont\setlength{\parskip}{0pt}\setlength{\aboverulesep}{1pt}\setlength{\belowrulesep}{1pt}
\R{2}{We are \CtwoConfidence{95\%} \CtwoConfidence{confident} that the population mean cycling-trip duration lies in $[\CtwoEndpoint{36.08},\CtwoEndpoint{43.92}]$ minutes, a range of plausible values under the model.}
\end{frame}
\section{Coverage evidence}
\begin{frame}[t]{Task 3 -- Assess the subsequent intervals}
\fontsize{10.5}{11.5}\selectfont\setlength{\parskip}{2pt} \At{8}{Coverage}{Five subsequent independent studies} sampled the \At{2}{PopulationMean}{same month's population}, each using 49 trips, the same standard deviation and \At{9}{Confidence}{95\% procedure}. A later comprehensive population investigation established an average of \At{2}{PopulationMean}{41 minutes}, enabling retrospective interval assessment.\par\smallskip
\begin{columns}[T,onlytextwidth]\begin{column}{.60\textwidth}\begin{tabular}{@{}rrrrl@{}}\toprule Study&$\bar{x}$ (min)&$L$ (min)&$U$ (min)&Check\\\midrule
\At{4}{Coverage}{1}&40&$\At{4}{Endpoint}{36.08}$&$\At{4}{Endpoint}{43.92}$&\R{4}{\CtwoCoverage{Cover}}\\
\At{5}{Coverage}{2}&42&$\At{5}{Endpoint}{38.08}$&$\At{5}{Endpoint}{45.92}$&\R{5}{\CtwoCoverage{Cover}}\\
\At{6}{Coverage}{3}&45&$\At{6}{Endpoint}{41.08}$&$\At{6}{Endpoint}{48.92}$&\R{6}{\CtwoCoverage{Miss}}\\
\At{7}{Coverage}{4}&38&$\At{7}{Endpoint}{34.08}$&$\At{7}{Endpoint}{41.92}$&\R{7}{\CtwoCoverage{Cover}}\\
\At{8}{Coverage}{5}&41&$\At{8}{Endpoint}{37.08}$&$\At{8}{Endpoint}{44.92}$&\R{8}{\CtwoCoverage{Cover}}\\
\bottomrule\end{tabular}\end{column}\begin{column}{.37\textwidth}\setlength{\parskip}{0pt}
\R{2}{Later established population mean: $\CtwoMean{\mu_{\mathrm{known}}}=\CtwoMean{41}$ minutes.}\par
\R{3}{Inclusive coverage condition: $\CtwoEndpoint{L_i}\leq\CtwoMean{41}\leq\CtwoEndpoint{U_i}$}.\par
\R{6}{Study 3 misses because $\CtwoEndpoint{L_3}=\CtwoEndpoint{41.08}>\CtwoMean{41}$.}
\end{column}\end{columns}
\begin{overlayarea}{\linewidth}{.8cm}
\only<4>{$\CtwoEndpoint{36.08}\leq\CtwoMean{41}\leq\CtwoEndpoint{43.92}$: study 1 covers.}
\only<5>{$\CtwoEndpoint{38.08}\leq\CtwoMean{41}\leq\CtwoEndpoint{45.92}$: study 2 covers.}
\only<7>{$\CtwoEndpoint{34.08}\leq\CtwoMean{41}\leq\CtwoEndpoint{41.92}$: study 4 covers.}
\only<8>{$\CtwoEndpoint{37.08}\leq\CtwoMean{41}\leq\CtwoEndpoint{44.92}$: study 5 covers.}
\R{9}{$\CtwoCoverCount(\CtwoConfidence{95\%})=\CtwoCoverage{4},\qquad\CtwoTotalCount(\CtwoConfidence{95\%})=\CtwoCoverage{5}$.}
\end{overlayarea}
\end{frame}
\begin{frame}[t]{Task 3 -- Observed coverage and absolute error}
\small
Four of the five subsequent intervals contain the later established fixed mean $\CtwoMean{\mu_{\mathrm{known}}}=\CtwoMean{41}$ minutes.
\[
\CtwoCoverCount(\CtwoConfidence{95\%})=\CtwoCoverage{4},\qquad\CtwoTotalCount(\CtwoConfidence{95\%})=\CtwoCoverage{5}.
\]
\begin{overlayarea}{\linewidth}{1.6cm}
\R{2}{$\CtwoCoverage{\mathrm{Cov}}(\CtwoConfidence{95\%})=\dfrac{\CtwoCoverage{4}}{\CtwoCoverage{5}}\times100\%$}\R{3}{$=\CtwoCoverage{80\%}$.}\par\medskip
\R{4}{$\CtwoError{E_{\rm abs}}(\CtwoConfidence{95\%})=|\CtwoCoverage{80}-\CtwoConfidence{95}|$}\R{5}{$=\CtwoError{15}$ \CtwoError{percentage points}.}
\end{overlayarea}
\R{6}{The \CtwoCoverage{observed} \CtwoCoverage{coverage} across five subsequent intervals is \CtwoCoverage{80\%}, \CtwoError{15 percentage points} below nominal \CtwoConfidence{95\%}. This finite-run discrepancy is consistent with the stated model.}

\end{frame}
\begin{frame}[t]{Task 3 -- What confidence means}
\small
\R{1}{\CtwoConfidence{Nominal} \CtwoConfidence{95\%} \CtwoConfidence{confidence} describes the procedure's long-run coverage: under the model, about 95\% of intervals from repeated fresh samples cover the fixed population mean.}\par\medskip
\R{2}{Each realized interval either contains that fixed parameter or misses it. Confidence is not a probability assigned to an already calculated interval.}\par\medskip
\R{3}{It is not the percentage of individual cycling-trip durations inside that interval.}\par\medskip
\R{4}{\CtwoCoverage{Finite-run} \CtwoCoverage{observed} \CtwoCoverage{coverage} need not equal nominal confidence or approach it monotonically as studies are added.}

\end{frame}
\section{Highlight guide}
\begin{frame}[t]{Key terms and highlight guide}
\BothLegends
\end{frame}
\end{document}
